\documentclass[11pt]{article}
\usepackage[utf8]{inputenc}
\usepackage{amsmath,amssymb,amsthm}
\usepackage{mathtools}
\usepackage{booktabs}
\usepackage{geometry}
\usepackage[round]{natbib}
\usepackage{hyperref}
\geometry{margin=1in}

\newtheorem{definition}{Definition}
\newtheorem{hypothesis}{Hypothesis}
\newtheorem{remark}{Remark}

\newcommand{\Spearman}{\rho_{s}}
\DeclareMathOperator{\AUROC}{AUROC}
\DeclareMathOperator{\sign}{sign}

\title{Convergent Validity of Localization Signals in Language-Model Inference}
\author{Kabir Singh Grewal}
\date{}

\begin{document}
\maketitle

%======================================================================
\begin{abstract}
A growing family of inference-time methods is justified by a claim of the
form \emph{``the information needed for this prediction lives at component
or token $X$.''} KV-cache eviction retains heavy-hitter tokens, head-aware
compression retains retrieval heads, and interpretability probes read
task-relevant content from specific positions or dictionary features. Each
method is validated against its own downstream objective; almost none are
validated against one another. We ask, \emph{per instance}, whether these
localization signals (i)~\textbf{agree} on where information lives,
(ii)~each \textbf{track a per-instance causal oracle} defined by token
masking---a causal analogue for the per-token information loss that eviction
induces---and (iii)~whether their
\textbf{disagreement carries information about KV-compression failure}
beyond trivial difficulty proxies such as context length, retrieval depth,
and model confidence. We organize signals into \emph{substrates} that answer
mechanistically distinct questions---input attribution (which context token
mattered) versus internal-confidence localization (where the model
represents that it knows)---and compute agreement only within a substrate.
Our central comparison is cross-family: whether attention-based and
dictionary-based (transcoder) signals, which process the model through
different machinery, converge on the same tokens, and whether either tracks
the causal oracle. Because distributed storage (high disagreement) is as
plausibly \emph{robust} to fixed-budget eviction as it is fragile, we treat
the direction of any disagreement--fragility relationship as an open
empirical question and pre-register a two-sided test. The contribution is
the per-instance convergent-validity framework and its causal grounding; any
predictive relationship between disagreement and compression failure is an
independent secondary result.
\end{abstract}

%======================================================================
\section{Introduction}

A localization claim asserts that the information a model uses for a specific
prediction is concentrated at an identifiable site---a set of tokens, a set
of attention heads, or a set of dictionary features. Such claims underwrite
a wide class of deployed inference methods. KV-cache eviction schemes retain
tokens judged to be heavy hitters and discard the rest; head-aware
compression retains heads identified as retrieval-relevant; and a parallel
line of interpretability work reads task content from probe directions or
sparse-dictionary features at specific positions. Each method is validated
against its own objective: an eviction scheme is judged by downstream task
accuracy under a memory budget, a probe by its detection accuracy. Almost
none are validated \emph{against each other}.

Two negative results bound what any \emph{single} localization method can
claim. \citet{hase2023} showed that the layer causal tracing identifies as
most important is largely not the layer at which model editing succeeds:
across most settings the localization signal and the intervention that works
are close to orthogonal, and even in the one variant where tracing is weakly
informative, it remains better to ignore the tracing result and edit an
early-to-mid-layer weight.
\citet{kantamneni2025} showed that sparse-autoencoder probes do not
reliably outperform logistic regression on raw activations across four
probing regimes (data scarcity, class imbalance, label noise, covariate
shift): the interpretable representation adds nothing over the raw signal
on that task. Both results bound a single method's localization claim.
Neither asks the meta-question this paper poses: if we line up the
localization signals practitioners actually use and ask them, on a specific
input, to point at the tokens or components that mattered, do they
\emph{agree with each other}, and does any of them agree with what a causal
intervention reveals?

We make this question precise as a problem of \emph{convergent validity}---a
construct-validation criterion borrowed from measurement theory: multiple
instruments intended to measure the same latent quantity should covary, and
should covary with an external criterion. Here the latent quantity is
``where the information for this prediction lives,'' the instruments are the
deployed localization signals, and the external criterion is a per-instance
causal oracle. We evaluate convergent validity \emph{per instance} rather
than in aggregate, on the \emph{deployed inference objects} rather than on
stylized circuit-discovery tasks, and with a causal oracle constructed as a
close analogue of the per-token information loss that eviction induces.

\paragraph{Contributions.}
\begin{enumerate}
  \item A per-instance convergent-validity framework for localization
  signals, with a substrate taxonomy that separates input attribution from
  internal-confidence localization and computes agreement only within a
  substrate (\S\ref{sec:formalization}).
  \item A causal-grounding criterion: a per-instance token-masking oracle
  constructed as a close analogue of the information loss KV eviction
  induces, against which each signal's fidelity is measured
  (\S\ref{sec:oracle}).
  \item A pre-registered, two-sided test of whether cross-method
  disagreement carries information about KV-compression failure beyond
  difficulty confounds, together with a battery of robustness checks that
  distinguish genuine localization plurality from dictionary
  reconstruction error (\S\ref{sec:fragility},~\S\ref{sec:robustness}).
\end{enumerate}

\paragraph{What we do not claim.}
We do not claim that any single method localizes well; \citet{hase2023} and
\citet{kantamneni2025} already bound this. We do not claim a direction for
the disagreement--fragility relationship; our test is two-sided by design.
We do not claim that all signals studied are deployed in production: KV
heavy-hitters and retrieval heads are deployment-grounded, whereas
dictionary features and confidence probes are candidate localization claims
drawn from the interpretability literature, included for the cross-substrate
comparison. And we do not treat within-family agreement as evidence on its
own, since signals computed from the same attention tensors are expected to
covary for partly mechanical reasons; the non-trivial comparison is
cross-family.

%======================================================================
\section{Problem Formalization}
\label{sec:formalization}

\paragraph{Instances.}
An instance $x$ consists of a rendered prompt with an identified set of
\emph{content tokens} $\mathcal{C}(x) = \{1,\dots,T\}$ (special and
chat-template tokens excluded by a fixed tokenizer-offset mask), a gold
answer span, and an answer string. We write $M$ for the language model and
$y(x)$ for the model's answer under the full KV cache.

\paragraph{Localization signals.}
A \emph{token-attribution signal} is a map producing a normalized importance
vector over content tokens,
\[
  s : x \;\longmapsto\; \mathbf{a}^{(s)}(x) \in [0,1]^{T},
\]
where larger $a^{(s)}_t$ means token $t$ is judged more important for
$y(x)$. A \emph{component signal} is a map producing a score over the query
heads of the model,
\[
  s : x \;\longmapsto\; \mathbf{h}^{(s)}(x) \in \mathbb{R}^{H},
\]
with $H$ the number of query heads. We study six signals. On the attention
side: accumulated attention (the heavy-hitter score that drives H2O-style
eviction) and two retrieval-head scores---the query-agnostic copy-paste
score of \citet{wu2024} and the query-focused QRscore of \citet{zhang2025},
which aggregates query-context attention and yields a distinct head set. On
the dictionary side: transcoder feature attribution, which natively produces
scores over \emph{features} and is reduced to a per-token vector by the
attribution defined in \S\ref{sec:formalization-reduction}. For the
internal-confidence question we include two signals---an
inference-time-intervention head set \citep{li2023iti} and an
exact-answer-token probe \citep{orgad2024}. Both were developed to localize
\emph{truthfulness} rather than ``knowing''; we treat them as candidate
proxies for internal-confidence localization, not as established measures of
it, and defend the relabeling in Remark~\ref{rem:confidence}.

\paragraph{Substrates.}
A \emph{substrate} is a set of signals that answer the same localization
question and are therefore commensurable. We define three:
\begin{itemize}
  \item \textbf{token-attribution}: which content token caused the answer
  (accumulated attention, retrieval-head--induced token importance from the
  Wu and QRHead head sets, and transcoder feature attribution reduced to
  tokens per \S\ref{sec:formalization-reduction});
  \item \textbf{component}: which heads carry the retrieval computation
  (the Wu head set, the QRHead head set, and the ITI head set);
  \item \textbf{answer-position}: where the model internally represents that
  it knows the answer (ITI directions, exact-answer-token probe).
\end{itemize}
Agreement is computed \emph{only within} a substrate; substrates are never
merged. The answer-position substrate addresses a distinct question
(internal confidence rather than input attribution) and is reported
separately.

\begin{remark}[The substrate distinction is a category boundary, not a
convenience]
Input attribution and internal-confidence localization are different
questions. ``Which context token caused the answer'' and ``where does the
model represent that it knows the answer'' need not have related answers:
a model may represent high confidence in a position that carries none of the
causal content. Merging them into a single ``where does the information
live'' score is a category error, and the separation is part of the
contribution.
\end{remark}

The ITI head set appears in two substrates---component and
answer-position---but is never double-counted within a single agreement
computation. It enters the component substrate as one instrument measuring
\emph{which heads carry retrieval}, and the answer-position substrate as one
instrument measuring \emph{where the answer is internally represented}.
These are two distinct construct measurements that happen to share an
instrument, which is admissible: agreement is always computed among the
members of one substrate, so ITI contributes to two separate agreement
matrices, not twice to the same one.

\begin{remark}[Truthfulness signals as candidate confidence proxies]
\label{rem:confidence}
ITI \citep{li2023iti} and the exact-answer-token probe \citep{orgad2024}
were developed to localize truthfulness, which is not identical to ``where
the model represents that it knows the answer.'' We therefore treat them as
candidate proxies rather than established measures. The strongest support for
the proxy is \citet{orgad2024}'s finding that a model can internally encode
the correct answer while generating an incorrect one: this dissociation
between internal representation and output is precisely the
``knows-but-does-not-say'' structure an internal-confidence signal should
capture, and it is what licenses reading the exact-answer-token probe as a
confidence localizer. We report the answer-position substrate separately
(RQ4) and do not let its interpretation bear on the token-attribution
results.
\end{remark}

\subsection{Reducing transcoder attribution to a per-token vector}
\label{sec:formalization-reduction}

The attention signals and the oracle are natively defined over content
tokens, but transcoder attribution is natively defined over \emph{features}:
it scores which sparse features contribute to the answer, and a feature may
fire at one position while encoding information gathered from many context
tokens. Placing transcoder attribution in the token-attribution substrate
therefore requires an explicit feature-to-token reduction, and the choice of
reduction is a load-bearing design decision rather than a formality, because
it determines what ``the transcoder localizes token $t$'' even means.

We adopt \emph{input-gradient attribution} as the primary reduction.
Let $\mathcal{F}(x)$ index the transcoder features active on instance $x$,
let $g_f(x)$ be the attribution of the answer to feature $f$ (its
contribution to the answer logit, e.g.\ via attribution patching on the
feature's activation), and let $z_f$ be that feature's pre-activation. The
contribution of feature $f$ is distributed over content tokens in proportion
to how much each token's embedding $\mathbf{e}_t$ drives the feature, measured
by the gradient of the pre-activation with respect to the embedding:
\[
  w_{f,t}(x)
  \;=\;
  \Bigl\lVert \frac{\partial z_f}{\partial \mathbf{e}_t} \Bigr\rVert_2 ,
  \qquad
  a^{(\mathrm{tc})}_t(x)
  \;=\;
  \sum_{f \in \mathcal{F}(x)}
  \lvert g_f(x) \rvert \;
  \frac{w_{f,t}(x)}{\sum_{t'} w_{f,t'}(x)} .
\]
The resulting vector $\mathbf{a}^{(\mathrm{tc})}(x)$ is min--max normalized to
$[0,1]^T$ like the other token-attribution signals. The intuition is that a
feature's answer-relevance $\lvert g_f \rvert$ is attributed back to the
content tokens it actually \emph{read}, traced through the gradient, rather
than to the single position where it happened to fire.

This is the choice the verification of Hypothesis~\ref{hyp:agreement}
depends on, so we treat it as a decision to validate, not to assume. As a
robustness check we re-run the cross-family agreement analysis under a
simpler \emph{firing-position} reduction, which attributes each feature's
$\lvert g_f \rvert$ to the position(s) at which it fires, and report whether
the agreement conclusion is stable across the two reductions. A conclusion
that flips between reductions would indicate that ``do attention and
transcoder converge on the same tokens'' is an artifact of the reduction
rather than a fact about the model; a stable conclusion licenses the
cross-family comparison.

Fix a substrate $S$ with member signals indexed $s \in S$. For each instance
$x$ and each pair $(s,s') \in S \times S$ we measure rank agreement between
the importance vectors restricted to content tokens (for token-attribution)
or to query heads (for component) using the Spearman rank correlation
\[
  \Spearman^{(s,s')}(x)
  \;=\;
  \operatorname{corr}\!\bigl(
    \operatorname{rank} \mathbf{a}^{(s)}(x),\,
    \operatorname{rank} \mathbf{a}^{(s')}(x)
  \bigr).
\]
We use Spearman as the primary metric because the signals are calibrated on
different scales and only their induced rankings are comparable. As
secondary metrics we report set overlap at top-$k$ via the Jaccard index
\[
  J_k^{(s,s')}(x)
  \;=\;
  \frac{\bigl| \mathrm{top}_k\,\mathbf{a}^{(s)}(x) \cap
               \mathrm{top}_k\,\mathbf{a}^{(s')}(x) \bigr|}
       {\bigl| \mathrm{top}_k\,\mathbf{a}^{(s)}(x) \cup
               \mathrm{top}_k\,\mathbf{a}^{(s')}(x) \bigr|},
\]
and rank-biased overlap (RBO) to weight agreement near the top of the
ranking, which is the regime a fixed-budget eviction scheme cares about.

\begin{definition}[Per-instance disagreement]
\label{def:disagreement}
The disagreement score of instance $x$ within substrate $S$ is one minus the
mean pairwise Spearman correlation over distinct signal pairs,
\[
  D_S(x)
  \;=\;
  1 \;-\;
  \frac{2}{|S|\,(|S|-1)}
  \sum_{\substack{s,s' \in S \\ s < s'}}
  \Spearman^{(s,s')}(x)
  \;\in\; [0,2].
\]
$D_S(x)$ is computed from the \emph{full-cache forward pass only}; it never
observes any compressed run. This is the structural leakage guarantee for
\S\ref{sec:fragility}.\footnote{We average the pairwise correlations
directly rather than via the Fisher $z$-transform. The $z$-transform would
remove a small bias in averaging correlations, but the resulting index is
monotone in the direct mean over the range we observe and the direct form is
the conventional choice for a disagreement score; we use it for
interpretability and report robustness to the $z$-transformed variant where
it matters.}
\end{definition}

\begin{remark}[Within-family agreement is partly mechanical]
Accumulated attention and retrieval-head token importance are both
functions of the same attention tensors; retrieval heads are identified by
their attention patterns. Their pairwise agreement is therefore expected to
be high for reasons that do not bear on localization plurality. We report it
as a controlled baseline. The informative quantity is the cross-family
correlation between an attention-derived signal and the transcoder signal,
which read the model through different mechanisms (attention routing versus
the MLP computation on the residual stream) and therefore localize
potentially different steps of the same task.
\end{remark}

To distinguish substantive agreement from chance under the finite token
budget, we calibrate each per-instance correlation against a permutation
null obtained by shuffling one ranking, yielding a per-instance $p$-value;
across instances we control the false discovery rate by the
Benjamini--Hochberg procedure, falling back to the Benjamini--Yekutieli
procedure if the per-instance tests show dependence that violates the
positive-dependence assumption.

%======================================================================
\section{The Causal Oracle}
\label{sec:oracle}

A localization signal can look authoritative on its own task yet fail to
identify the tokens that causally determine the answer. We therefore define
a per-instance causal criterion against which every signal is scored. We
construct the criterion to be as close as possible to the information loss a
KV-eviction scheme induces, while being explicit that it is an analogue and
not the identical operation (see the discussion of mechanical differences
below).

\begin{definition}[Causal effect under masking]
\label{def:oracle}
For instance $x$ and content token $t$, let $\ell\bigl(y(x) \mid x\bigr)$ be
the model's log-probability of its full-cache answer, and let $x^{\setminus t}$
be the forward pass in which token $t$'s contribution is \emph{masked}: its
key/value entries are removed from attention at every layer, so that no query
at any position can attend to $t$, while all other tokens' cached
representations are left in place. The causal effect of $t$ is the induced
drop in answer log-probability,
\[
  E(t)
  \;=\;
  \ell\bigl(y(x) \mid x\bigr)
  \;-\;
  \ell\bigl(y(x) \mid x^{\setminus t}\bigr).
\]
\end{definition}

We use masking rather than Gaussian noise on activations, which would push
activations off-distribution \citep{zhang2024}; masking removes a token's
availability without injecting off-manifold signal. Masking is also the
operation by which H2O validated its heavy-hitter criterion, which aligns the
oracle with the deployment method it targets. As a robustness check we report
stability of the resulting ranking under symmetric token replacement
\citep{zhang2024}, an alternative on-manifold corruption that swaps the input
token rather than masking its cache entry; conclusions are reported only where
they hold across both corruptions.

\paragraph{The oracle is an analogue of eviction, not the same operation.}
Masking a token's key/value entries is the closest single-token analogue of
KV eviction available within a clean causal-effect definition, but it is not
identical to what an eviction scheme does, and we do not claim it is. Three
mechanical differences remain. First, eviction \emph{drops} a stored
key/value vector, whereas even masking leaves the question of what was already
computed at earlier positions untouched in a different way than a deployed
cache eviction interacts with the generation loop; the oracle reads a single
counterfactual forward pass, not the autoregressive trajectory an eviction
scheme actually runs under. Second, eviction operates per-layer on cached
$K/V$ as decoding proceeds, removing availability to \emph{future} queries,
whereas the oracle's masking is applied uniformly across the prompt's forward
pass. Third, the symmetric-swap robustness variant changes the residual
stream of neighboring positions through their attention to the swapped token,
propagating from the embedding upward, which eviction does not do at all.
Because the causal-grounding contribution rests on the oracle, we therefore
read $E(t)$ as an \emph{upper-bound proxy} for the per-token information loss
eviction induces, and we test signals against it on that understanding rather
than treating fidelity to $E$ as fidelity to eviction itself. The empirical
bridge to eviction is the flip test of \S\ref{sec:fragility}, which uses the
deployed schemes directly.

\paragraph{Fidelity.}
Each token-attribution signal $s$ is scored by the rank correlation between
its importance vector and the oracle effect,
\[
  F(s; x)
  \;=\;
  \Spearman\!\bigl( \mathbf{a}^{(s)}(x),\, \mathbf{E}(x) \bigr),
\]
where $\mathbf{E}(x) = (E(1),\dots,E(T))$. Computing $E(t)$ exactly requires
one masked forward pass per content token. We therefore admit an
attribution-patching estimator $\widehat{E}(t)$---a first-order Taylor
approximation requiring two forward passes and a single backward pass,
independent of the number of tokens \citep{nanda2022,kramar2024}---for the
full sweep, gated by a fidelity requirement: $\widehat{E}$ is used only where
\[
  \Spearman\!\bigl( \widehat{\mathbf{E}}(x),\, \mathbf{E}(x) \bigr) \;\ge\; 0.8
\]
on a verification subsample; otherwise we fall back to exact masking on the
adjudication subset. This gate guarantees the criterion that adjudicates all
downstream claims is itself trustworthy, and directly addresses
attribution patching's known failure mode of underestimating effects for
distributed, high-magnitude activations \citep{kramar2024}.

\begin{remark}[The oracle is what adjudicates convergent validity]
Pairwise agreement between two signals does not by itself establish that
either is correct. Two signals can agree because both are right (both track
$\mathbf{E}$) or because both share the same artifact. Fidelity to the
oracle separates these cases. Convergent validity in the strong sense we
test requires both \emph{covariation among signals} and \emph{covariation of
each signal with the external causal criterion}.
\end{remark}

%======================================================================
\section{Disagreement and Compression Fragility}
\label{sec:fragility}

The deployment-facing question is whether the per-instance disagreement
$D_S(x)$, measured before any compression, carries information about whether
a fixed-budget KV-eviction scheme will flip a correct answer to an incorrect
one.

\begin{definition}[Flip]
For a compression method $m$ (e.g.\ H2O, SnapKV) at memory budget $b$, let
$c_{\mathrm{full}}(x) \in \{0,1\}$ indicate whether the full-cache answer is
correct and $c_{m,b}(x) \in \{0,1\}$ whether the compressed answer is
correct, both under a deterministic correctness metric (token-level $F_1
\ge 0.5$ as primary, exact match as robustness). Instance $x$
\emph{flips} under $(m,b)$ if
\[
  \mathrm{flip}_{m,b}(x) \;=\; \mathbf{1}\!\left[
    c_{\mathrm{full}}(x) = 1 \;\wedge\; c_{m,b}(x) = 0
  \right].
\]
\end{definition}

Flips are defined only on instances answered correctly under the full cache,
so that the event is a genuine compression-induced failure rather than a
pre-existing error. We model the flip probability with logistic regression
and ask whether $D_S(x)$ adds explanatory power over a confound-only model:
\begin{align}
  \text{(confounds only)} \quad
  &\operatorname{logit} \Pr[\mathrm{flip}_{m,b}(x)=1]
   = \beta_0 + \boldsymbol{\beta}_z^\top \mathbf{z}(x), \\[4pt]
  \text{(confounds + disagreement)} \quad
  &\operatorname{logit} \Pr[\mathrm{flip}_{m,b}(x)=1]
   = \beta_0 + \boldsymbol{\beta}_z^\top \mathbf{z}(x) + \gamma\,D_S(x),
\end{align}
where $\mathbf{z}(x)$ collects the difficulty proxies (context length,
retrieval depth, full-cache answer confidence). The test of interest is on
$\gamma$, evaluated by a likelihood-ratio test between the two nested models
and by the change in cross-validated predictive performance ($5$-fold,
held-out $\AUROC$). We report the confound-only $\AUROC$ \emph{before}
introducing $D_S(x)$, so that any incremental contribution is visible and
not absorbed silently. We report both the LRT and the held-out $\AUROC$
change deliberately: a coefficient can be statistically significant under the
LRT while adding little held-out discrimination, and the two answers need not
agree, so we treat a significant $\gamma$ as necessary but not sufficient for
a claim of practical predictive value.

\begin{remark}[Why the sign is left free]
A high $D_S(x)$ admits two opposing mechanistic readings. If disagreement
reflects that the answer-relevant information is \emph{distributed} across
many tokens or components, a fixed-budget eviction is \emph{less} likely to
remove all of it, predicting robustness ($\gamma < 0$). If disagreement
reflects that the signals are jointly \emph{mis-localized}, whatever the
scheme retains is more likely to miss the load-bearing tokens, predicting
fragility ($\gamma > 0$). We do not have an a priori reason to prefer either,
so the hypothesis below is two-sided and pre-registered as such.
\end{remark}

%======================================================================
\section{Hypotheses}
\label{sec:hypotheses}

We state the claims the experiments will adjudicate. None is assumed; each is
a hypothesis whose resolution---in either direction---is reportable.

\begin{hypothesis}[Within-substrate agreement structure]
\label{hyp:agreement}
Within the token-attribution substrate, cross-family pairwise agreement
(attention-derived vs.\ transcoder) is systematically lower than
within-family agreement (attention-derived vs.\ attention-derived), and
varies with task difficulty and context length. Formally, the mean
cross-family Spearman is bounded away below from the mean within-family
Spearman.
\end{hypothesis}

\begin{hypothesis}[Causal grounding]
\label{hyp:fidelity}
Each signal's fidelity $F(s;x)$ to the oracle is the quantity that decides
whether observed (dis)agreement reflects genuine localization plurality.
Three mutually exclusive outcomes are admissible and each is reportable:
(a)~the dictionary signal tracks the oracle and an attention signal does not,
(b)~the reverse, or (c)~neither tracks the oracle strongly. The data select
among these.
\end{hypothesis}

\begin{hypothesis}[Disagreement carries fragility information]
\label{hyp:fragility}
The coefficient $\gamma$ on $D_S(x)$ in the flip model is nonzero after
controlling for difficulty confounds:
\[
  H_0\!:\ \gamma = 0
  \qquad\text{vs.}\qquad
  H_1\!:\ \gamma \ne 0 .
\]
The test is two-sided. A failure to reject is itself informative: it bounds
what cross-method disagreement adds beyond difficulty proxies. A degenerate
case---insufficient variance in $D_S(x)$ to test---is reported as such rather
than as a null, and in that case the primary result reverts to
Hypothesis~\ref{hyp:fidelity}.
\end{hypothesis}

%======================================================================
\section{Robustness: Distinguishing Plurality from Reconstruction Error}
\label{sec:robustness}

The central threat to any cross-family disagreement result is that the
transcoder signal disagrees with attention \emph{because the transcoder is a
lossy approximation of the residual stream}, not because localization is
genuinely plural. We address this with tests built into the design rather
than by argument.

\paragraph{(i) Oracle adjudication.}
Disagreement is evidence of plurality only if the disagreeing signals do not
both collapse onto the oracle. By Hypothesis~\ref{hyp:fidelity}, a signal
enters the analysis as informative only if it tracks $\mathbf{E}(x)$; a
transcoder that fails the fidelity gate is dropped and reported as a finding,
not treated as a source of disagreement.

\paragraph{(ii) Reconstruction-error correlation.}
The lossy hypothesis makes a falsifiable prediction: disagreement should
concentrate on instances the transcoder reconstructs worst. Let
$r(x)$ be the transcoder's reconstruction error on the answer-relevant
positions. If
\[
  \Spearman\!\bigl( r(x),\, D_S(x) \bigr)
\]
is large and positive, the lossy reading is supported; if it is weak, the
disagreement persists on cleanly reconstructed instances and the lossy
reading is undercut. This test is evidence against the lossy reading, not
proof against it: if $r(x)$ has low variance it cannot correlate with
anything, and if lossiness degrades the ranking only at non-answer positions
while $r(x)$ is measured at answer-relevant positions, the two can decouple
even when lossiness does drive disagreement. We therefore treat the oracle
adjudication in~(i) and the high-fidelity check in~(iii) as the primary
adjudicators, with this correlation as corroborating evidence.

\paragraph{(iii) High-fidelity feature check.}
On an instance where a specific transcoder feature reconstructs the
answer-relevant positions with near-zero error, persistent disagreement with
the attention signal cannot be attributed to reconstruction loss. A single
such case is a qualitative refutation of the lossy reading.

\paragraph{Captured variance.}
Throughout, we report the answer-relevant variance the transcoder captures,
as a necessary (not sufficient) precondition; (i)--(iii) are what carry the
argument.

%======================================================================
\section*{Acknowledgments}
[omitted in draft]

\bibliographystyle{plainnat}
% \bibliography{lcv}  % uncomment once .bib is assembled
\begin{thebibliography}{99}
\bibitem[Hase et al.(2023)]{hase2023}
P.~Hase, M.~Bansal, B.~Kim, and A.~Ghandeharioun.
\newblock Does localization inform editing? Surprising differences in
causality-based localization vs.\ knowledge editing in language models.
\newblock NeurIPS 2023. arXiv:2301.04213.

\bibitem[Kantamneni et al.(2025)]{kantamneni2025}
S.~Kantamneni, J.~Engels, S.~Rajamanoharan, M.~Tegmark, and N.~Nanda.
\newblock Are sparse autoencoders useful? A case study in sparse probing.
\newblock ICML 2025. arXiv:2502.16681.

\bibitem[Wu et al.(2025)]{wu2024}
W.~Wu et al.
\newblock Retrieval head mechanistically explains long-context factuality.
\newblock ICLR 2025. arXiv:2404.15574.

\bibitem[Zhang et al.(2025)]{zhang2025}
W.~Zhang, F.~Yin, H.~Yen, D.~Chen, and X.~Ye.
\newblock Query-focused retrieval heads improve long-context reasoning and
re-ranking.
\newblock EMNLP 2025. arXiv:2506.09944.

\bibitem[Zhang et al.(2023)]{h2o}
Z.~Zhang et al.
\newblock H\textsubscript{2}O: Heavy-hitter oracle for efficient generative
inference of large language models.
\newblock NeurIPS 2023. arXiv:2306.14048.

\bibitem[Li et al.(2024)]{snapkv}
Y.~Li et al.
\newblock SnapKV: LLM knows what you are looking for before generation.
\newblock arXiv:2404.14469.

\bibitem[Dunefsky et al.(2024)]{transcoders}
J.~Dunefsky, P.~Chlenski, and N.~Nanda.
\newblock Transcoders find interpretable LLM feature circuits.
\newblock NeurIPS 2024. arXiv:2406.11944.

\bibitem[Li et al.(2023)]{li2023iti}
K.~Li, O.~Patel, F.~Vi\'egas, H.~Pfister, and M.~Wattenberg.
\newblock Inference-time intervention: Eliciting truthful answers from a
language model.
\newblock NeurIPS 2023. arXiv:2306.03341.

\bibitem[Orgad et al.(2025)]{orgad2024}
H.~Orgad et al.
\newblock LLMs know more than they show: On the intrinsic representation of
LLM hallucinations.
\newblock ICLR 2025. arXiv:2410.02707.

\bibitem[Nanda(2022)]{nanda2022}
N.~Nanda.
\newblock Attribution patching: Activation patching at industrial scale.
\newblock Technical report, 2022.

\bibitem[Kram\'ar et al.(2024)]{kramar2024}
J.~Kram\'ar, T.~Lieberum, R.~Shah, and N.~Nanda.
\newblock AtP*: An efficient and scalable method for localizing LLM behaviour
to components.
\newblock arXiv:2403.00745.

\bibitem[Zhang \& Nanda(2024)]{zhang2024}
F.~Zhang and N.~Nanda.
\newblock Towards best practices of activation patching in language models:
Metrics and methods.
\newblock ICLR 2024. arXiv:2309.16042.
\end{thebibliography}

\end{document}
